\documentclass[a4paper,11pt,dvipdfmx]{jsarticle}
\usepackage{amsmath, amssymb, amsthm}
\usepackage{bm}
\usepackage{mathtools}
\usepackage{comment}
\usepackage{algorithm}
\usepackage{algorithmic}
\usepackage{bm}
\usepackage{mathtools}
\usepackage{enumitem}
\mathtoolsset{showonlyrefs}
\usepackage[dvipdfmx, colorlinks=true, linkcolor=blue, citecolor=blue, urlcolor=blue, setpagesize=false]{hyperref}
\usepackage{pxjahyper}
% --- 定理・定義環境 ---
\theoremstyle{definition}
\newtheorem{definition}{定義}[section]
\newtheorem{theorem}[definition]{定理}
\newtheorem{lemma}[definition]{補題}
\newtheorem{assumption}{仮定}
\newtheorem{proposition}[definition]{命題}
\newtheorem{requirement}[definition]{要請}
\newtheorem{example}[definition]{例}
\newcommand{\fU}{\underline{f}}
\newcommand{\gU}{\underline{g}}
\newcommand{\uU}{\underline{u}}
\newcommand{\aU}{\underline{a}}
\newcommand{\bU}{\underline{b}}
\newcommand{\cU}{\underline{c}}
\newcommand{\dU}{\underline{d}}
\newcommand{\vU}{\underline{v}}
\newcommand{\zeroU}{\underline{0}}
\newcommand{\C}[2]{C_{{#1},{#2}}}
\newcommand{\cy}[3]{\mathcal{C}_{{#1},{#2},{#3}}}
\newcommand{\inv}{^{-1}}
\newcommand{\HHX}{\hat{H}_X}
\newcommand{\HHZ}{\hat{H}_Z}
\newcommand{\tr}{^\top}

% =========================================
\begin{document}

\title{ゼミ記録ノート}
\author{}
\date{\today}
\maketitle


\section{3/13}
\subsection{進捗}
\begin{itemize}
  \item 条件ABの一般解を求めた。
  \item 置換群全体での群的構成を行ったが、可換性条件とガースの両立が難しい。
  \item APMの不動点に着目した可換群を考えてみたが、これだと可換性条件を実現できない。
  \item 一つの置換行列を2つの領域に分け、2つのAPMで1つのAPMを構成することを考えたが、$f,g$が見つかっていない。
\end{itemize}

\subsection{今後のタスク}
\begin{itemize}
  \item 条件ABを満たす解が存在するための$\aU$の条件を考える。
  \item 現状条件CはABを満たすものを生成後、検査するしかないが、よく発生する閉サイクルがあるなら、これを条件ABと同じ段階で満たすようにする。
  \item 空間結合符号について考える
\end{itemize}

\subsection{今後の方針}
\begin{itemize}
  \item 条件ABを満たす解は高速で見つけられるのでそこからILPで条件Cを満たすものを求める。
  \item 条件Cの一部を解を求める段階に入れる
  \item 空間結合符号について調べる
\end{itemize}



\section{2/27}
\subsection{進捗}
\begin{itemize}
  \item 一つの特殊解からほかの解を見つけられるようになった。
  \item ガースを大きくする方法を考えた。
\end{itemize}


\subsection{今後のタスク}
\begin{itemize}
  \item $\aU$をランダムに生成し、それに対して整数計画法で特殊解を求め、性能を示すものを探す。また、UTCBC以外の長さ8のサイクルを含まないものが存在するのか考える。
  \item ガースが10になるような符号の構成法を考える。今の案としては2つ。1つは置換行列を配置せず、より疎な行列にする。2つ目は、アクティブ行列の取り出し方(実質的に置換行列の配置)を変える。
  \item 置換の関数として2次の多項式を考える。
\end{itemize}

\subsection{今後の方針}
\paragraph{目指すべき符号}
サイズが小さく、ガースが大きく、最小距離が大きい符号が望ましい。
サイズの小ささは、実用性に繋がる。ベース行列が大きくなっても、$P$を小さくできればサイズは結果的に小さくなる。
\paragraph{$\fU,\gU$の一般解の探索}
特殊解から見つけた他の解は、グラフ同型になり、性能が同じになってしまう。

逐次的な探索アルゴリズムを実装してみる。

二次の置換は、群として閉じていない。よって、一般的な置換群を扱うことになる。
今は線型方程式を使って解こうとしているが、群の性質を使って解く方法もあるのではないか。
\paragraph{ガースを大きくする方法}
疎なベース行列：岡田が、列重みに、多元化を用いた構成を研究している。\\
多元化をしない方法、空間結合符号(帯行列)の性能がいいらしい\\
列重み2と3を混ぜる\\

\paragraph{他の方針}
最小距離の下界の証明

\section{2/4}
\subsection{進捗}
\begin{itemize}
  \item \cite{kasai2025quantum}ではガースが12の符号について、UTCBC以外の長さ12のサイクルを含まない、という条件も課していた、\cite{kasai2026breakingorthogonalitybarrierquantum}ではガース8で、長さ4と6のサイクルを含まない、と記述されている。つまり、UTCBC以外の長さ8のサイクルは存在してよい、という認識で正しいのか。
  \item 整数計画法は特殊解を求めるためのアルゴリズムなので、一般解を求めるという今回のゴールには向いてなさそう。スミス標準形を使って一般解を求め、制約条件に基づき係数の条件を求める。
\end{itemize}

\subsection{今後のタスク}
\begin{itemize}
  \item 整数計画法で特殊解を求める。
  \item スミス標準形を勉強。
  \item サイクルのカウントが正しいのか確かめる(長さ8のサイクルがすべてUTCBCで正しいのか？)。
  \item 性能評価のやり方を勉強する。
\end{itemize}

\subsection{今後の方針}
\begin{itemize}
  \item 長さ6のサイクルがあるのではないか：関数$f(x)$と行列$F$の関係について、\cite{kasai2025quantum}とは行と列が逆、アクティブ部でのみサイクルを考える。
  \item UTCBC以外の長さ8のサイクルを避けるべきか？：今は長さ6以外のサイクルを含まないことを条件にしているが、もしUTCBC以外のサイクルも避ければ、性能が良くなるかも？(ならないかも)、エラーがどのサイクル由来かを調べられるとよい
  \item 整数計画法で求めた特殊解はランダムネスが小さい
  \item APMを使っている時点でランダムネスは減少しているのでAPMに縛らないとランダムネスを増やせる
  \item 他の改良案：Pを小さく(2,3以外の素因数を使う)、ガースを大きく(アクティブ行列の取り出し方を変える。)
\end{itemize}


\end{document}